% Private replacement body for the current B.24--B.26; no public source edit.
\subsection{A full-state completion for every Alice angle}
\label{sec:c1-c141-all-alice}

\begin{theorem}\label{thm:c1-c141-all-tests}
Retain the original stationary C14.1 source and all its intersections:
$d=1$, $\lambda_{\rm ev}=n=0$, the native strict gaps,
$m,C,A,\alpha>0$, $D,\beta\ge0$, and
\[
 x=k^2,\qquad z=(x\zeta)^{-1},\qquad\sigma=\eta^{-1}
          \quad\hbox{positive source infinitesimals}.
\]
There is no restriction on $H=h/\sigma$. The complete homogeneous
parent is a finite tensor sum of squares at every independent
projective testing pair and every real or complex state. The literal
faces $D=0$, $\beta=0$ and all testing and state-coordinate zeros
are included.
\end{theorem}

\begin{proof}
We first treat the corner in which the Bob half-angle is projectively
unbounded. Put $w=1/v$, including $w=0$ at $v=\infty$.
Thus $w$ is infinitesimal or zero. If $u$ is infinitesimal or zero,
Theorem~\ref{thm:c1-unification-bounded-z-all-H} applies directly.
For every other Alice test, $p=1/u$ is bounded; $p=0$ includes
$u=\infty$. The exact angle gauges allow $p,w\ge0$: changing the
sign of $p$ together with $Y,Z$, or changing the sign of $w$ together
with $X,Z$, preserves the homogeneous polynomial below.

\paragraph{The whole polynomial in the inverse-Alice chart.}
Write $D_p=1+p^2$ and $D_w=1+w^2$. The homogeneous parent is
\[
\begin{split}
 \widetilde F={}&4m(pW-pwX+Y-wZ)^2+4CD_w(W-pY)^2
                         +4AD_pD_wY^2\\
 &+(4D+2)(wW+X-pwY-pZ)^2
 +(4\alpha+2)D_p(Y-wZ)^2\\
 &+(4\beta+2)D_w(pX+Z)^2-2D_pD_wZ^2.
\end{split}                                                    \tag{IA1}
\]
For $p\ne0$ this is $p^2$ times the reciprocal parent
\eqref{eq:c1-reciprocal-parent}; the displayed polynomial defines its
literal projective endpoint at $p=0$. Set
\[
 q=\sqrt z,\quad Q=1-b,\quad r=bz,\quad H_0=1+zQ,
\]
\[
 E=qW,\quad F=Y/q,\quad G=wZ/q,\quad
 R=X+wW,\quad S=pX,\quad J=F-bG.
\]
Only pure-source inverses occur in these roots. Direct expansion of
the complete signed part of \eqref{eq:c1-reciprocal-parent} gives
\[
 K=R^2+D_wS^2+z(G-F)^2+p^2D_wzF^2
                    -2pE(G+w^2F)+2qwS(G-F).             \tag{IA2}
\]
Thus \eqref{eq:c1-reciprocal-parent}, in this chart, is $2K$ plus
the six original positive event terms. This identity keeps all four
state coordinates.

\paragraph{Three leading event weights complete the signed core.}
Define the positive source coefficients
\[
 C_0=\frac{z}{2(1-r)},\qquad
 A_0=\frac{H_0}{2bz},\qquad
 \alpha_0=\frac1{2Qz},\qquad B=\frac{H_0}{bQ}.
\]
The opening-cleared estimates (BE9) imply
\[
 C/C_0\sim1,\qquad A/A_0\sim1,\qquad\alpha/\alpha_0\sim1.
\]
In particular the pure-source number
\[
 f=\min\{1,C/C_0,A/A_0,\alpha/\alpha_0\}
\]
is positive, $f\sim1$, and each original event weight is at least
$f$ times its corresponding leading weight. This uses neither a
bound on $H$ nor positivity of $D$ or $\beta$.

The decisive cancellation is the following exact identity:
\[
\begin{split}
 &2z(G-F)^2+2p^2zF^2-4pEG
 +\frac2{1-r}(E-pzF)^2\\
 &\hspace{8mm}+\frac{2H_0}{b}(1+p^2)F^2
                    +\frac2Q(1+p^2)(F-G)^2\\
 &=\frac2{1-r}(E-pzJ-pG)^2
                  +2B(1+p^2)J^2+2H_0G^2.
\end{split}                                                    \tag{IA3}
\]
It follows by substituting $F=J+bG$ and $r=bz$; no contact or state
equation is used. In particular it joins the negative product to
three whole event squares before any coefficient is estimated.

Put
\[
 J_* =\sqrt B J,\quad G_* =\sqrt{H_0}G,\quad
 a=\frac z{\sqrt B},\quad c=\frac1{\sqrt{H_0}},\quad
 L=aJ_*+cG_*=zJ+G.
\]
The native inequality $r<1$ gives the exact contraction
\[
 1-a^2-c^2=\frac{zQ(1-r)}{H_0}\ge0.                  \tag{IA4}
\]
Consequently the source coefficients $a,c$ are bounded, including
arbitrarily small $b$ or $Q$.

\paragraph{An exact positive remainder and the entire error.}
For clarity define
\[
\begin{split}
 \mathcal X={}&4m(pW-pwX+Y-wZ)^2
       +4D(wW+X-pwY-pZ)^2+4\beta D_w(pX+Z)^2\\
 &+4(C-fC_0)D_w(W-pY)^2
       +4(A-fA_0)D_pD_wY^2\\
 &+4(\alpha-f\alpha_0)D_p(Y-wZ)^2,
\end{split}
\]
and retain the nonnegative expression
\[
\begin{split}
 \Pi={}&\mathcal X
 +2(1-f)\{R^2+S^2+z(G-F)^2+p^2zF^2\}\\
 &+2w^2S^2+2p^2w^2zF^2
 +4fC_0w^2(W-pY)^2+4fA_0D_pw^2Y^2.
\end{split}                                                    \tag{IA5}
\]
All its source coefficients are nonnegative; $D_p,D_w$ and the
remaining factors separate in the testing field. Each term is
therefore a tensor sum of squares. Equations (IA1)--(IA3) give
the full equality
\[
\begin{split}
 \widetilde F={}&f\mathcal Q+\Pi+\mathcal E,\\
 \mathcal Q={}&2R^2+2S^2+\frac2{1-r}(E-pL)^2
                         +2(1+p^2)J_*^2+2G_*^2,\\
 \mathcal E={}&-4(1-f)pEG-4pw^2EF+4qwS(G-F).
\end{split}                                                    \tag{IA6}
\]
In particular the last state coupling has the original sign from
(IA2). Nothing has been restricted to a state subspace.

\paragraph{A uniform reserve for every bounded inverse angle.}
Choose a real constant $P\ge1$ with $p\le P$ and put
\[
 t=p/P,\qquad E_*=E/P,\qquad
 g=(R,S,E_*,J_*,G_*)^{\mathsf T}.
\]
Then $0\le t\le1$. Define
\[
 \mathcal B=2R^2+2S^2+2(E_*-tL)^2
                         +2(1+t^2)J_*^2+2G_*^2.
\]
The difference $\mathcal Q-\mathcal B$ is exactly
\[
 2\left(\frac{P^2}{1-r}-1\right)(E_*-tL)^2
                    +2(P^2-1)t^2J_*^2\in\Sigma T^2.
                                                               \tag{IA7}
\]
Moreover, the contraction (IA4) gives the explicit reserve
\[
\begin{split}
 \mathcal B-\tfrac12g^{\mathsf T}g
 ={}&\tfrac32(R^2+S^2)
       +\tfrac32(E_*-4tL/3)^2+\tfrac32J_*^2+\tfrac56G_*^2\\
 &+\tfrac23(1-t^2)G_*^2+\tfrac43t^2J_*^2\\
 &+\tfrac23t^2\{(1-a^2-c^2)(J_*^2+G_*^2)
                                   +(cJ_*-aG_*)^2\}.
\end{split}                                                    \tag{IA8}
\]
Every radical needed on the right is in one labelled field.
In particular no square root of a mixed coefficient is taken.

It remains to put the complete error in the same five roots.
Its only nonzero upper-triangle matrix entries are
\[
\begin{array}{c|c}
 (E_*,G_*)&-2(1-f)P^2t/\sqrt{H_0}
                         -2P^2btw^2/\sqrt{H_0}\\
 (E_*,J_*)&-2P^2tw^2/\sqrt B\\
 (S,J_*)&-2qw/\sqrt B\\
 (S,G_*)& 2qQw/\sqrt{H_0}.
\end{array}                                                    \tag{IA9}
\]
The pure-source coefficient matrices at $t,tw^2,w$ have absolute
row bounds $2(1-f)P^2$, $4P^2$, and $4$, respectively. Here use
$1/\sqrt B\le1/2$, $b/\sqrt{H_0}\le1$, and
\[
 q^2/B=rQ/H_0\le1,\qquad q^2Q^2/H_0\le1.
\]
The first bound is a source infinitesimal. In the other two terms,
$tw^2$ and $w$ are testing infinitesimals. Thus (IA9) defines a
bounded tensor matrix with zero real residue, without imposing an
order on the tensor product or comparing source and testing rates.
The remaining Gram matrix is $fI/2$ plus this error; its residue is
$I/2$. The strict bounded-Gram lemma proves its tensor sum of
squares. Restoring (IA5), (IA7) and (IA8) proves the whole corner
for every homogeneous real state. The roots themselves need not
be bounded or independent.

\paragraph{The remaining Bob angles and the endpoints.}
On C14.1, (BE9) and $z$ infinitesimal imply that $A$ and $\alpha$
are unbounded. The adjacent-event theorem transported by
(CE9)--(CE10) in Section~\ref{component:p89fd3edb3edd} therefore
covers every Bob test that is not projectively unbounded, with all
Alice tests. This and the corner just proved exhaust the projective
testing pairs.

No inverse of $D$, $\beta$ or a state coordinate occurs. The
identities specialize literally at $p=0$, $w=0$ and zero state
coordinates; $u=0$ was covered separately at the start. Real and
imaginary state parts give complex states. Finally restore the
positive pure-field homogenization factors as in the original
parent normalization. For unit states the inverse-Alice angular
factor is $1/(D_pD_w)$, which remains positive at $p=0$; arbitrary
nonzero homogeneous states have the additional pure-testing state
norm factor. This proves the stated full C14.1 domain.
\end{proof}
